\documentclass[11pt]{article}

\usepackage[T1]{fontenc}
\usepackage[utf8]{inputenc}
\usepackage{lmodern}
\usepackage{microtype}
\usepackage{amsmath,amssymb,amsthm}
\usepackage{mathtools}
\usepackage{geometry}
\usepackage{hyperref}
\usepackage{enumitem}

\geometry{margin=1.05in}
\hypersetup{
  colorlinks=true,
  linkcolor=black,
  citecolor=black,
  urlcolor=blue,
  pdftitle={A Balayage Counterexample to Nikliborc's Problem 150}
}

\newtheorem{theorem}{Theorem}[section]
\newtheorem{remark}[theorem]{Remark}

\title{\textbf{A Balayage Counterexample to Nikliborc's Problem 150}}
\author{}
\date{September 2026}

\newcommand{\publicationstatus}{%
\begin{center}
\fbox{\begin{minipage}{0.92\linewidth}
\small
\textbf{AI EXPLORATORY MATHEMATICAL NOTE}\\[2pt]
\textbf{Status:} E1 - AI cross-checked; \textbf{NO INDEPENDENT HUMAN REVIEW}\\
\textbf{Version:} v1.0\\
\textbf{First public release:} PENDING\\
\textbf{Related Lean formalization:} a specific off-center-sphere counterexample is kernel-checked relative to an explicit balayage axiom. The broader smooth-domain construction in this note is not fully formalized; see the end matter.\\
\end{minipage}}
\end{center}
\medskip
}

\hypersetup{pdfauthor={}}

\begin{document}

\maketitle

\publicationstatus

\begin{abstract}
Nikliborc's Problem 150 in the \emph{Scottish Book} asks whether reflection symmetry, near infinity, of the Newtonian single-layer potential of a closed surface forces the supporting surface and the surface density to have the same reflection symmetry. We construct a counterexample.
\end{abstract}

\begin{quote}
\small
\textbf{Feedback.} This note has not undergone independent human mathematical review. Readers who know relevant prior literature, identify an error, or wish to check the argument are especially encouraged to contact the coordinators listed at the end of the paper or to open an issue in the repository.
\end{quote}


\section{The printed problem}

Nikliborc's Problem 150 is printed on p.~250 of the second edition of the \emph{Scottish Book} \cite[p.~250]{ScottishBook}. It starts with a closed surface $S$ and a \emph{continuous} function $f$ on $S$, and considers the single-layer potential
\[
  V(P)=\int_S \frac{f(M)}{r_{PM}}\,d\sigma_M.
\]
The hypothesis is that whenever two points $P_1,P_2$ lie outside a sufficiently large sphere and are symmetric with respect to a plane $\pi$, one has $V(P_1)=V(P_2)$. The problem asks one to prove that $\pi$ is a plane of symmetry of $S$ and that $f$ takes equal values at symmetric points of $S$.

Thus the printed problem assumes only continuity of $f$; no positivity hypothesis is imposed. The counterexample below is stronger in this respect, because its density is smooth and strictly positive.

\section{The counterexample}

Let $R_\pi$ denote reflection in a plane $\pi\subset\mathbb R^3$.

\begin{theorem}\label{thm:main}
There exist a smooth closed surface $S\subset\mathbb R^3$, a smooth strictly positive density
\[
  f:S\longrightarrow (0,\infty),
\]
and a plane $\pi$ such that the single-layer potential
\[
  V(x)=\int_S \frac{f(y)}{|x-y|}\,d\sigma(y)
\]
satisfies
\[
  V(R_\pi x)=V(x)
\]
for every sufficiently large $|x|$, whereas
\[
  R_\pi(S)\neq S.
\]
Consequently, this construction is a counterexample to the first asserted conclusion of Nikliborc's Problem 150.
\end{theorem}

\begin{proof}
Choose a plane $\pi$ through the origin. Let $D\subset\mathbb R^3$ be a bounded domain with smooth boundary such that
\[
  0\in D,
\]
but
\[
  R_\pi(\partial D)\neq \partial D.
\]
For instance, one may take a sufficiently small nonsymmetric smooth radial perturbation of the unit ball. Put
\[
  S=\partial D.
\]

Sweep the unit point mass $\delta_0$ at the origin onto $S$ by Newtonian balayage. The resulting measure $\mu$ is supported on $S$, has total mass one, and has the same Newtonian potential as $\delta_0$ at exterior points:
\begin{equation}\label{eq:balayage}
  U^\mu(x)
  =\int_S \frac{d\mu(y)}{|x-y|}
  =\frac{1}{|x|},
  \qquad x\in\mathbb R^3\setminus\overline D.
\end{equation}

For completeness, let us record the standard boundary density underlying this balayage. Normalize the Dirichlet Green function $G_D(x,y)$ by
\[
  -\Delta_y G_D(x,y)=\delta_x,
  \qquad
  G_D(x,y)\sim \frac{1}{4\pi|x-y|}
  \quad (y\to x),
\]
with $G_D(x,y)=0$ for $y\in\partial D$. If $n_y$ denotes the outward unit normal, then the balayage measure of $\delta_0$ has density
\begin{equation}\label{eq:density}
  d\mu(y)=f(y)\,d\sigma(y),
  \qquad
  f(y)=-\frac{\partial G_D}{\partial n_y}(0,y).
\end{equation}
For a smooth boundary, the standard boundary regularity of the Green function gives a smooth density $f$. Moreover,
\[
  G_D(0,y)>0\quad (y\in D\setminus\{0\}),
  \qquad
  G_D(0,y)=0\quad (y\in\partial D),
\]
so the Hopf boundary point lemma, applied with the outward normal convention, yields
\[
  \frac{\partial G_D}{\partial n_y}(0,y)<0
  \qquad (y\in\partial D).
\]
Hence
\[
  f(y)>0\qquad (y\in S).
\]
This is stronger than the continuity required in the printed problem.

Combining \eqref{eq:balayage} and \eqref{eq:density},
\begin{equation}\label{eq:surfacepotential}
  V(x)
  =\int_S \frac{f(y)}{|x-y|}\,d\sigma(y)
  =\frac{1}{|x|},
  \qquad x\in\mathbb R^3\setminus\overline D.
\end{equation}

Because $\pi$ passes through the origin,
\[
  |R_\pi x|=|x|.
\]
Choose $R>0$ so large that
\[
  D\cup R_\pi(D)\subset B_R(0).
\]
If $|x|>R$, then both $x$ and $R_\pi x$ lie in $\mathbb R^3\setminus\overline D$. Therefore \eqref{eq:surfacepotential} gives
\[
  V(R_\pi x)
  =\frac{1}{|R_\pi x|}
  =\frac{1}{|x|}
  =V(x).
\]
Thus $V$ has the reflection symmetry required by the hypothesis of Problem 150 outside a sufficiently large sphere.

On the other hand,
\[
  R_\pi(S)=R_\pi(\partial D)\neq\partial D=S.
\]
Therefore $\pi$ is not a plane of symmetry of the supporting surface. The first requested conclusion of Problem 150 already fails, so no separate failure of the density-symmetry conclusion is needed.
\end{proof}

\section{Remarks on the mechanism}

The counterexample is a standard manifestation of nonuniqueness for inverse exterior-potential problems. Equation \eqref{eq:surfacepotential} says that the exterior field of the surface distribution is exactly the field of a unit point mass at the origin. Hence, on every sufficiently distant region, it is invariant not merely under the chosen reflection but under the full orthogonal group $O(3)$, although the supporting surface may be chosen without the corresponding symmetry.

The restriction to sufficiently large $|x|$ is essential in the formulation of the symmetry statement. Since $D$ need not be invariant under $R_\pi$, an exterior point $x$ can have $R_\pi x\in D$. Thus \eqref{eq:surfacepotential} does not by itself imply $V(R_\pi x)=V(x)$ for every exterior $x$. The argument above proves exactly what the printed problem assumes: symmetry outside one sufficiently large sphere.

\section{Historical provenance}

The potential-theoretic mechanism is classical. MacMillan develops Poincar\'e's method of balayage in Chapter VI of \emph{The Theory of the Potential}; see in particular the discussion beginning with ``The Balayage of a Sphere'' on pp.~313--318 \cite[pp.~313--318]{MacMillan}. More directly, Problem~1 on p.~323 asks the reader to show by balayage that on any given closed surface there is a unique unit-mass distribution that is centrobaric with respect to a prescribed interior point \cite[p.~323, Problem~1]{MacMillan}. This is precisely the classical exterior-field mechanism used above.

Accordingly, Theorem~\ref{thm:main} is not presented as a new theorem in potential theory. The observation here is that this classical construction contradicts the literal conclusion of Nikliborc's Problem 150. We have not established whether this particular application to Problem 150 was previously recorded, so no priority claim is made for this counterexample.


\section*{Acknowledgments}
Chuyang Chen and Chenhao Bian are grateful to their advisor, Fei Yu, for his guidance on this project.

\section*{Independent review and Lean verification notice}
\begingroup
\small
\begin{center}
\begin{minipage}{0.94\linewidth}
\centering
\textbf{\large No independent human mathematical review has been performed.}
\end{minipage}
\end{center}

\medskip
\noindent The accompanying Lean file compiles under Lean 4.33.1 / Mathlib version v4.33.1 (exact mathlib revision \texttt{0df444a360eaa60ab8c11dca51a86af692955474}).

\medskip
\noindent The accompanying Lean file formalizes an off-center-sphere counterexample relative to the explicit axiom \texttt{Nikliborc150.classical\_balayage\_on\_sphere}; the broader smooth-domain mechanism in this note is not fully formalized.

\medskip
\noindent\textbf{Successful Lean verification establishes correctness of the formalized statement relative to its assumptions. It does not by itself establish that the formalized statement is exactly equivalent to the historical problem as originally formulated.}
\endgroup

\section*{Provenance and contacts}
\begingroup\small\sloppy
\begin{description}
\item[Project curator.] Fei Yu.
\item[AI-generation coordinators and contacts.] Chuyang Chen (\href{mailto:22535018@zju.edu.cn}{22535018@zju.edu.cn}); Chenhao Bian (\href{mailto:bch26@mails.tsinghua.edu.cn}{bch26@mails.tsinghua.edu.cn}).
\item[Mathematical draft generation.] Mathematical drafts were generated with GPT-5.6 Sol under their supervision.
\item[Lean code generation.] The accompanying Lean source was generated and revised with GPT-5.6 Sol under their supervision. This is a code-generation provenance statement and is separate from mathematical review.
\item[Lean build/kernel execution.] The local build and kernel-check commands reported below were executed by Chuyang Chen in the recorded project environment.
\item[Human mathematical verification.] NONE.
\item[Statement-correspondence audit.] NONE. No independent human audit has certified that the natural-language statement and the formal Lean statement are exactly equivalent.
\end{description}
\medskip
\noindent Comments, corrections, historical references, and independent verification are very welcome. For long-term correspondence, readers are encouraged to use the repository issue tracker as well as the contacts above.
\endgroup

\section*{Lean formal verification record}

\begingroup\small\sloppy
This record separates the natural-language statement, the formal Lean statement, kernel checking of the proof term, and the separate audit of the correspondence between the first two. Kernel acceptance is not treated as human verification of the original problem.

\begin{description}
\item[Formalization status.] Conditional F1 relative to one explicit external mathematical axiom.
\item[Lean source.] \texttt{\detokenize{PaperFormalization/Scottish 150/Main.lean}}.
\item[Principal formal theorems.] \mbox{}\\
\texttt{\detokenize{Nikliborc150.counterexample}}\\
\texttt{\detokenize{Nikliborc150.printed_problem_150_is_false}}.
\item[Build command.] \texttt{\detokenize{lake env lean "PaperFormalization/Scottish 150/Main.lean"}}.
\item[Build/kernel result.] PASS.
\item[Lean version.] Lean 4.33.1 (commit 819816b2e0a3bf405af45ae5c7af2491d8f5bee6).
\item[Library snapshot.] mathlib v4.33.1, exact revision 0df444a360eaa60ab8c11dca51a86af692955474.
\item[Project commit.] PENDING -- the final verified working tree has not yet been committed.
\item[Kernel axiom audit.] Both principal theorems depend on \texttt{propext}, \texttt{Classical.choice}, \texttt{Quot.sound}, and the user-defined mathematical axiom \texttt{\detokenize{Nikliborc150.classical_balayage_on_sphere}}.
\item[External axiom statement.] For a Euclidean sphere of radius $r>0$ centered at $c$ and containing the origin, the axiom supplies a continuous strictly positive density on the sphere whose single-layer Newtonian potential equals $\|x\|^{-1}$ outside the corresponding closed ball.
\item[Scope note.] Everything after this explicit balayage interface is proved in Lean. This is therefore a kernel-checked derivation relative to the stated external axiom, not an axiom-free Lean proof of the balayage theorem itself. The Lean counterexample uses an off-center Euclidean sphere, while the paper presents the broader smooth-domain balayage mechanism.
\end{description}

\noindent\textbf{Interpretation of the label.} F1 applies to the stated conditional formal result. It does not certify the external balayage axiom, and F2 would additionally require a human statement-correspondence audit.
\endgroup

\begin{thebibliography}{9}

\bibitem{ScottishBook}
R. D. Mauldin (ed.),
\emph{The Scottish Book: Mathematics from the Scottish Caf\'e, with Selected Problems from the New Scottish Book},
2nd ed., Birkh\"auser/Springer, Cham, 2015.
DOI: \href{https://doi.org/10.1007/978-3-319-22897-6}{10.1007/978-3-319-22897-6}.

\bibitem{MacMillan}
W. D. MacMillan,
\emph{The Theory of the Potential},
McGraw--Hill, New York and London, 1930.
See pp.~313--318 for Poincar\'e's method of balayage and p.~323, Problem~1, for the unit-mass construction on an arbitrary closed surface.

\end{thebibliography}

\end{document}
